\section*{Chapitre \ref{ch:b1:filters-transfer-functions} --- Filtres et fonctions de transfert}

\begin{solution}{exo:b1:filters-transfer-functions:1}
$0.5 \to \qty{-6.0}{dB}$~; $10 \to \qty{20}{dB}$~; $1/\sqrt2 \to \qty{-3.0}{dB}$~;
$100 \to \qty{40}{dB}$. $\qty{-40}{dB} \to 0.01$~; $\qty{+6}{dB} \to 2.0$.
Cascade~: $-6 - 20 = \qty{-26}{dB}$, soit un rapport $0.05$.
\end{solution}

\begin{solution}{exo:b1:filters-transfer-functions:2}
$f_c = 1/(2\pi RC) = 1/(2\pi \times 10^4 \times 1.6\times10^{-8}) =
\qty{1.0}{kHz}$. \omterm{def:b1:filters-transfer-functions:H}{Gains} $-10\log(1 + (f/f_c)^2)$~: \qty{100}{Hz}~:
\qty{-0.04}{dB}~; \qty{1}{kHz}~: \qty{-3}{dB}~; \qty{10}{kHz}~: \qty{-20}{dB}~;
\qty{100}{kHz}~: \qty{-40}{dB}.
\end{solution}

\begin{solution}{exo:b1:filters-transfer-functions:3}
(a) passe-bas ($C$ court-circuite la sortie à haute $f$)~; (b) passe-haut
($C$ s'ouvre à basse $f$)~; (c) passe-bas du second ordre ($L$ s'ouvre à
haute $f$, $C$ court-circuite)~; (d) passe-bande ($C$ bloque les basses
$f$, $L$ les hautes $f$).
\end{solution}

\begin{solution}{exo:b1:filters-transfer-functions:4}
$4E/n\pi$~: \qty{6.37}{V} à \qty{1}{kHz}, \qty{2.12}{V} à \qty{3}{kHz},
\qty{1.27}{V} à \qty{5}{kHz}. Le \omterm{def:b1:wave-propagation:signal}{signal} carré est antisymétrique par
décalage d'une demi-période ($s(t + T/2) = -s(t)$), ce que les
harmoniques paires ne sont pas.
\end{solution}

\begin{solution}{exo:b1:filters-transfer-functions:5}
À \qty{100}{Hz} et à ses harmoniques utiles, $\omega \ll \omega_c$~:
$\underline H \approx j\omega/\omega_c$, une dérivée divisée par
$\omega_c$. La pente du triangle vaut $\pm 4Ef = \pm\qty{400}{V/s}$, donc
la sortie est un \omterm{def:b1:wave-propagation:signal}{signal} carré d'\omterm{def:b1:wave-propagation:signal}{amplitude} $400/(2\pi \times 10^4) =
\qty{6.4}{mV}$.
\end{solution}

\begin{solution}{exo:b1:filters-transfer-functions:6}
Comme $\omega \gg \omega_c$~: $u_s \approx (1/\tau)\int u_e\,\dd t$, avec $\tau =
1/(2\pi \times 10) = \qty{16}{ms}$. L'intégrale de $\pm E$ est un
triangle qui monte de $ET/2$ par demi-période~: \omterm{def:b1:wave-propagation:signal}{amplitude} crête à crête
$ET/2\tau$, donc \omterm{def:b1:wave-propagation:signal}{amplitude}
$ET/4\tau = 1.0 \times 10^{-3}/(4 \times 0.016) = \qty{16}{mV}$ autour de
zéro.
\end{solution}

\begin{solution}{exo:b1:filters-transfer-functions:7}
$f_0 = \qty{5.03}{kHz}$. $Q = \sqrt{L/C}/R = 1/\sqrt2$~: $R = \sqrt2 \times
316 = \qty{447}{\ohm}$. En $f_0$~: \qty{-3}{dB}~; en $10f_0$~: $1/\sqrt{1 +
10^4} \to \qty{-40}{dB}$.
\end{solution}

\begin{solution}{exo:b1:filters-transfer-functions:8}
$G = 1/\sqrt{1 + Q^2(x - 1/x)^2}$ avec $x = f/3$~: à \qty{1}{kHz}, $x = 1/3$,
$Q(x - 1/x) = -53$, $G = 0.019$~; à \qty{3}{kHz}, $G = 1$~; à \qty{5}{kHz},
$x = 5/3$, $21$, $G = 0.047$. Sortie~: \qty{2.1}{V} à \qty{3}{kHz}, avec
\qty{0.12}{V} à \qty{1}{kHz} et \qty{0.06}{V} à \qty{5}{kHz} --- une
sinusoïde à \qty{3}{kHz} presque pure.
\end{solution}

\begin{solution}{exo:b1:filters-transfer-functions:9}
Aux nœuds~: avec $a = jRC\omega$, le premier condensateur voit $R$ en
amont et $R + 1/jC\omega$ en aval~: $\underline H = 1/(1 + 3a + a^2)$, soit
$1/[1 - (\omega/\omega_c)^2 + 3j\omega/\omega_c]$. En $\omega_c$~: $1/3j$,
$G = 1/3$, \qty{-9.5}{dB}~; le produit des deux \omterm{def:b1:filters-transfer-functions:H}{gains} séparés vaudrait
$1/2$, soit \qty{-6}{dB}. $f_c = \qty{1.0}{kHz}$.
\end{solution}

\begin{solution}{exo:b1:filters-transfer-functions:10}
Sous $\omega_1$~: \qty{0}{dB}~; entre $\omega_1$ et $\omega_2$~: montée à
\qty{20}{dB} par décade~; au-dessus de $\omega_2$~: plateau à $20\log(\omega_2/\omega_1) =
\qty{40}{dB}$. Phase $\arctan(\omega/\omega_1) - \arctan(\omega/\omega_2)$~:
en $\omega = 10^3$, $\arctan 10 - \arctan 0.1 = 84^\circ - 6^\circ =
78^\circ$.
\end{solution}

\begin{solution}{exo:b1:filters-transfer-functions:11}
Bien au-dessus de la coupure, $G \approx f_c/f$~; le taux d'ondulation
après filtrage vaut
\[
\frac{(4E/3\pi)(f_c/100)}{2E/\pi} = \frac{2}{3}\,\frac{f_c}{100} < 0.01,
\]
d'où $f_c < \qty{1.5}{Hz}$ et $\tau = 1/2\pi f_c > \qty{0.1}{s}$. La charge elle-même joue le rôle de
$R$~: un condensateur à ses bornes ne coûte ni chute de tension ni
puissance, là où une résistance en série gaspillerait les deux.
\end{solution}

\begin{solution}{exo:b1:filters-transfer-functions:12}
$f_c = 1/(2\pi RC) = 1/(2\pi \times 1.6\times10^5 \times 10^{-7}) =
\qty{10}{Hz}$, $\tau = RC = \qty{16}{ms}$. Pendant un plateau, le
condensateur se charge à travers $R$~: la sortie décroît en $\eu^{-t/\tau}$
depuis chaque front. Affaissement sur $T/2$~: $1 - \eu^{-T/2\tau}$~; à \qty{50}{Hz}
($T/2 = \qty{10}{ms}$)~: $1 - \eu^{-0.63} = 47\%$~; à \qty{1}{kHz}
($T/2 = \qty{0.5}{ms}$)~: $3\%$.
\end{solution}

\begin{solution}{pb:b1:filters-transfer-functions:1}
\textbf{1.} $\underline H = \underline U_{\text{HP}}/\underline U_e$~;
$G_{\mathrm{dB}} = 20\log\abs{\underline H}$.

\textbf{2.} $\underline H_L = R/(R + jL\omega) = 1/(1 + j\omega/\omega_c)$,
$\omega_c = R/L$.

\textbf{3.} $L = R/\omega_c = 8.0/(2\pi \times 2000) = \qty{0.64}{mH}$.

\textbf{4.} $\underline H_H = R/(R + 1/jC\omega) = jRC\omega/(1 + jRC\omega)$,
$\omega_c = 1/RC$~; $C = 1/(R\omega_c) = \qty{9.9}{\micro F} \approx
\qty{10}{\micro F}$.

\textbf{5.} En $f_c$~: \qty{-3}{dB} chacun. À \qty{200}{Hz} ($x = 0.1$)~:
graves \qty{-0.04}{dB}, aigus \qty{-20}{dB}. À \qty{20}{kHz} ($x = 10$)~:
graves \qty{-20}{dB}, aigus \qty{-0.04}{dB}.

\textbf{6.} $G_L^2 + G_H^2 = (1 + x^2)/(1 + x^2) = 1$~: à charges égales,
les deux puissances totalisent celle qu'absorberait un unique
haut-parleur de \qty{8}{\ohm} --- partagée, non perdue.

\textbf{7.} $-45^\circ$ (graves) et $+45^\circ$ (aigus)~: $90^\circ$
d'écart~; en $f_c$, les deux membranes ne bougent pas ensemble.

\textbf{8.} \qty{\pm 6}{dB} par octave (\qty{20}{dB} par décade).

\textbf{9.} Diviseur entre $jL\omega$ et $R \parallel 1/jC\omega =
R/(1 + jRC\omega)$~: $\underline H = R/[R + jL\omega(1 + jRC\omega)] =
1/(1 - LC\omega^2 + jL\omega/R)$.

\textbf{10.} $\omega_0 = 1/\sqrt{LC}$~; $L\omega/R = x/Q$ donne $Q =
R/L\omega_0 = R\sqrt{C/L}$.

\textbf{11.} $\abs{\underline H}^{-2} = (1 - x^2)^2 + x^2/Q^2 = 1 + x^4 +
x^2(1/Q^2 - 2) = 1 + x^4$ pour $Q^2 = \tfrac12$~; en $x = 1$, $G = 1/\sqrt2$.

\textbf{12.} $L = R/(Q\omega_0) = 8\sqrt2/(1.257\times10^4) = \qty{0.90}{mH}$~;
$C = 1/(L\omega_0^2) = \qty{7.0}{\micro F}$.

\textbf{13.} $x = 2$~: $1/\sqrt{17}$, \qty{-12.3}{dB}~; $x = 4$~: $1/\sqrt{257}$,
\qty{-24.1}{dB}~: soit \qty{12}{dB} par octave.

\textbf{14.} Échanger $L$ et $C$ échange $x \leftrightarrow 1/x$ (et
$Q$ reste)~: $\underline H_H = -x^2/(1 - x^2 + jx/Q)$.

\textbf{15.} En $x = 1$~: $\underline H_L = Q/j = -jQ$ ($-90^\circ$),
$\underline H_H = jQ$ ($+90^\circ$)~: les deux haut-parleurs sont en
opposition de phase à la coupure et s'annuleraient acoustiquement~;
l'aigu est donc câblé en inversant sa polarité.

\textbf{16.} Les harmoniques au-dessus de \qty{2}{kHz}~: $n \geq 10$
(\qty{2.2}{kHz} et au-delà) partent surtout vers l'aigu.

\textbf{17.} $9 \times 220 = \qty{1980}{Hz} \approx f_c$~: partagée à
parts égales, \qty{-3}{dB} de chaque côté.

\textbf{18.} Graves ($G_L$)~: $x = 0.5$~: $0.89$~; $x = 1.5$~: $0.55$~; $x = 2.5$~:
$0.37$. Aigus ($G_H$)~: $0.45$, $0.83$, $0.93$.

\textbf{19.} L'aigu~: son condensateur série bloque le continu. Le
haut-parleur de graves encaisse la composante continue (\qty{1}{V} sur
\qty{8}{\ohm}~: \qty{125}{mA}, \qty{0.125}{W} de chaleur et une membrane
décentrée).

\textbf{20.} Premier ordre, $x = 0.25$~: $G_H = 0.25/\sqrt{1.0625} = 0.24$
(\qty{-12}{dB}), soit $6\%$ de la puissance. Second ordre~: $0.0625/\sqrt{1 +
0.0039} = 0.06$ (\qty{-24}{dB}), $0.4\%$. Un haut-parleur d'aigus donné
pour quelques watts reçoit $6\%$ d'une note grave jouée fort avec le
\omterm{def:b1:filters-transfer-functions:H}{filtre} du premier ordre --- il peut brûler.

\textbf{21.} $\underline Z = 8 + jL_{vc}\omega$~: à \qty{2}{kHz}, $8 + 6.3j$,
$\abs{\underline Z} = \qty{10.2}{\ohm}$~; à \qty{8}{kHz}, $8 + 25.1j$,
\qty{26.3}{\ohm}.

\textbf{22.} $\underline H = \underline Z/(\underline Z + jL\omega)$~: à
\qty{2}{kHz}, $10.2/\abs{8 + 14.3j} = 10.2/16.4 = 0.62$ (\qty{-4.2}{dB})~;
à \qty{8}{kHz}, $26.3/\abs{8 + 57.3j} = 0.45$ (\qty{-6.9}{dB}) au lieu de
\qty{-12.3}{dB}~: l'\omterm{def:b1:sinusoidal-impedance:impedance}{impédance} croissante de la bobine mobile contrarie
la bobine série, et la pente s'aplatit à quelques dB par octave.

\textbf{23.} \omterm{def:b1:sinusoidal-impedance:impedance}{Admittances}~: $1/(R + jL_{vc}\omega) + jC_z\omega/(1 + jRC_z\omega)$~;
avec $C_z = L_{vc}/R^2$, le second terme vaut $jL_{vc}\omega/[R(R +
jL_{vc}\omega)]$, et la somme fait $(R + jL_{vc}\omega)/[R(R + jL_{vc}\omega)]
= 1/R$. $C_z = 0.50\times10^{-3}/64 = \qty{7.8}{\micro F}$.

\textbf{24.} $U = \sqrt{PR} = \sqrt{400} = \qty{20}{V}$ efficaces. En
$f_c$, chaque haut-parleur en reçoit la moitié~: \qty{25}{W}. À
\qty{500}{Hz}, l'aigu reçoit $6\%$~: \qty{3}{W}.

\textbf{25.} Premier ordre~: deux composants, puissances
complémentaires, pentes douces de \qty{6}{dB/oct}, $90^\circ$ entre les
haut-parleurs, et des graves qui fuient vers l'aigu. Second ordre~:
quatre composants, \qty{12}{dB/oct}, $180^\circ$ à la coupure (il faut
inverser l'aigu), protection bien meilleure. Dans les deux cas,
l'inductance du haut-parleur de graves doit être domptée par un réseau
de Zobel pour que le dimensionnement tienne.
\end{solution}
